\documentclass{article}

\PassOptionsToPackage{numbers}{natbib}
\usepackage[preprint]{neurips_2026}

\usepackage[utf8]{inputenc}
\usepackage[T1]{fontenc}
\usepackage{hyperref}
\usepackage{url}
\usepackage{booktabs}
\usepackage{amsfonts}
\usepackage{amsmath}
\usepackage{amssymb}
\usepackage{nicefrac}
\usepackage{microtype}
\usepackage{xcolor}
\usepackage{graphicx}
\usepackage{enumitem}
\usepackage{listings}

\lstset{
  basicstyle=\small\ttfamily,
  breaklines=true,
  frame=single,
  xleftmargin=1em,
}

\title{Convergence Criteria and Accuracy Metrics\\
  for GNN AC Power-Flow Surrogates}

\author{Changhun Kim}

\begin{document}
\maketitle

\begin{abstract}
  This note documents (1) how Newton--Raphson (NR) solvers define
  convergence --- verified against pandapower v3.2.1 --- (2) how other
  solver families define it, (3) why a single metric is insufficient for
  a GNN surrogate, and (4) the recommended reporting table.  Special
  attention is given to the numerical conditioning issues that arise in
  stiff, multi-voltage grids such as LVN.
\end{abstract}

%==============================================================================
\section{How pandapower (and the NR Family) Define Convergence}
%==============================================================================

\subsection{The Criterion}

pandapower's convergence check (v3.2.1, verified on source) is:

\begin{lstlisting}[language=Python]
def _check_for_convergence(F, tol):
    return linalg.norm(F, inf) < tol      # default tol = 1e-8 MVA
\end{lstlisting}

Convergence is declared when the infinity-norm of the \emph{masked} power
mismatch vector falls below the tolerance:

\begin{equation}
  \|F\|_\infty < \texttt{tol},
  \qquad
  F =
  \begin{bmatrix}
    \Delta P_i & \text{(PV + PQ buses)} \\
    \Delta Q_i & \text{(PQ buses only)}
  \end{bmatrix}
  \label{eq:nrconv}
\end{equation}

with $\Delta P_i = P_i^{\rm set} - \mathrm{Re}\{V_i \overline{(YV)_i}\}$
and $\Delta Q_i = Q_i^{\rm set} - \mathrm{Im}\{V_i \overline{(YV)_i}\}$.

Three properties of \eqref{eq:nrconv} are critical:

\paragraph{1. The mismatch is masked.}
Slack bus $P$ and $Q$, and PV bus $Q$, are \emph{excluded} from $F$.
These are free variables in the AC-PF problem; they adjust to whatever
the network requires and cannot be in violation.  Including them produces
spurious large residuals (the slack bus absorbs all system imbalance and
can show $|\Delta P| \sim 30$\,p.u.\ even at a perfectly converged
solution).

\paragraph{2. The norm is $\ell^\infty$ (worst single bus).}
A single badly modelled bus can prevent convergence even if all others
are accurate.  This is strict but physically meaningful for transmission
grids.

\paragraph{3. The tolerance is research-grade tight.}
The default $\texttt{tol} = 10^{-8}$\,MVA is far tighter than what
industry solvers typically require (Section~\ref{sec:other}).  All
$V_{\rm newton}$ reference solutions in the training data satisfy this
criterion.

\subsection{Equivalent per-unit statement}

If $S_{\rm base} = 100\,\mathrm{MVA}$, the criterion becomes
$\|F_{\rm pu}\|_\infty < 10^{-10}$\,p.u.  For LVN with
$S_{\rm base} = 1\,\mathrm{MVA}$, the same physical tolerance corresponds
to $10^{-8}$\,p.u.\ --- i.e.\ the per-unit numbers are 100$\times$ larger
than HV case numbers for the same physical accuracy.

%==============================================================================
\section{Convergence Criteria Across Solver Families}
\label{sec:other}
%==============================================================================

\begin{table}[h]
\centering
\caption{Convergence criteria used by different AC power-flow solvers.}
\label{tab:solvers}
\begin{tabular}{p{3.8cm}p{3.2cm}p{2.2cm}p{2.8cm}}
\toprule
Solver & Criterion & Tolerance & Notes \\
\midrule
pandapower / PYPOWER / MATPOWER
  & $\|[\Delta P, \Delta Q]\|_\infty$
  & $10^{-8}$ pu/MVA
  & Masked; research-grade \\[3pt]
PSS/E, PowerWorld
  & Max mismatch [MW, MVAr]
  & $0.1$--$1$ MVA
  & Engineering tolerance \\[3pt]
OpenDSS, GridLAB-D
  & $\|\Delta V\|_\infty$
  & $10^{-4}$ pu
  & Fixed-point / sweep \\
\bottomrule
\end{tabular}
\end{table}

Two philosophies dominate:

\begin{itemize}[noitemsep]
  \item \textbf{Power-mismatch norm} (NR / transmission): convergence means
    KCL is satisfied at every constrained bus.  This is the same quantity
    used in the PIGNN physics loss.
  \item \textbf{Voltage-increment norm} (fixed-point / distribution):
    convergence means voltages are no longer changing between iterations.
    This does not directly imply KCL satisfaction.
\end{itemize}

Because the reference solutions $V_{\rm newton}$ in this project were
generated by pandapower (NR), the appropriate benchmark is the
power-mismatch criterion.

%==============================================================================
\section{Why One Metric Is Not Enough for a GNN Surrogate}
%==============================================================================

The LVN experiments (Section~7 of the results document) demonstrated
concretely that \textbf{prediction accuracy and physical feasibility can
diverge}:

\begin{itemize}[noitemsep]
  \item Pure MSE training: $\theta$ RMSE $= 0.79^\circ$ (good accuracy),
    physics loss $\gg 1$ (bad feasibility).
  \item Combined loss: $\theta$ RMSE $\approx 2.1^\circ$ (worse accuracy),
    physics loss $\approx 0.9$ (better feasibility).
\end{itemize}

A surrogate that is accurate but not feasible copies NR output without
independently satisfying KCL --- useful as a warm-start initialiser but
not as a stand-alone solver.  A surrogate that is feasible but not
accurate satisfies KCL at a different operating point than the NR
reference.

\textbf{Both axes must be reported.}

\subsection{Axis 1: Prediction Accuracy (vs.\ NR Reference)}

\begin{equation}
  |V|_{\rm err} = \hat{V} - V_{\rm newton},
  \qquad
  \theta_{\rm err} = \hat{\theta} - \theta_{\rm newton}
\end{equation}

Report: MAE and max for both $|V|$ [pu] and $\theta$ [deg].

This axis answers: \emph{did the surrogate find the NR solution?}

\subsection{Axis 2: Physical Feasibility (KCL Satisfaction)}

Compute the masked mismatch from the surrogate's predicted voltages
$(\hat{V}, \hat{\theta})$:

\begin{align}
  V_c &= \hat{V} \cdot e^{j\hat{\theta}}, \quad
  S_c  = V_c \odot \overline{Y V_c} \\
  \Delta P_i &= P_i^{\rm set} - \mathrm{Re}\{S_{c,i}\}
    \quad \text{(PV + PQ buses)} \\
  \Delta Q_i &= Q_i^{\rm set} - \mathrm{Im}\{S_{c,i}\}
    \quad \text{(PQ buses only)}
\end{align}

Report: mean, 95th percentile, and max of $|\Delta P|$ and $|\Delta Q|$
in both MW/MVAr and p.u.

This axis answers: \emph{would a standard NR solver call this converged?}

\subsection{The Most Interpretable Single Number: Pseudo-Convergence Rate}

Report the \textbf{fraction of test cases (or buses) whose masked mismatch
falls below an engineering tolerance}, e.g.\ $1\,\mathrm{MVA}$:

\begin{equation}
  \text{convergence rate}
  = \frac{|\{i : |\Delta P_i| < \epsilon \;\text{and}\; |\Delta Q_i| < \epsilon\}|}
         {|\{\text{constrained buses}\}|}
  \label{eq:convrate}
\end{equation}

with $\epsilon = 10^{-3}$\,p.u.\ ($0.1\,\mathrm{MVA}$ at 100\,MVA base).

\medskip
\noindent Example: \emph{``92\% of PQ buses satisfy the 1\,MVA engineering
tolerance.''} This mirrors how industry solvers judge a solution and is
far more informative than a raw $\ell^\infty$ norm.

%==============================================================================
\section{Critical Caveats for Stiff Multi-Voltage Grids (LVN)}
%==============================================================================

\subsection{Never Report $\ell^\infty$ Alone}

In the LVN Heo1 grid, the maximum diagonal admittance entry reaches
$|Y_{ii}| \approx 5.6 \times 10^6$\,p.u.  The power mismatch at a bus
with voltage error $\epsilon$ is approximately:

\begin{equation}
  |\Delta S_i| \approx |Y_{ii}| \cdot |\epsilon|
\end{equation}

A small voltage error of $0.7^\circ$ can produce a mismatch of
$\sim 10^5$\,p.u.\ at the single stiffest bus, even when the
mean mismatch is $\sim 4 \times 10^{-4}$\,p.u.  Reporting only the
$\ell^\infty$ norm hides this discrepancy.

\textbf{Always report: mean/median, 95th percentile, and max.}

\subsection{Always Apply the Correct Bus Mask}

Omitting the slack/PV mask produces a spurious large $\ell^\infty$
value at the slack bus even for a perfectly converged solution.  The
masking must match pandapower's convention:

\begin{itemize}[noitemsep]
  \item $\Delta P$: zero the \textbf{slack bus} only.
  \item $\Delta Q$: zero the \textbf{slack bus and all PV buses}.
\end{itemize}

\subsection{Use Normalized Residuals for Loss and Feasibility Metrics}

For training on stiff multi-voltage grids, the raw $\ell^2$ mismatch
loss is ill-conditioned because $|Y_{ii}|$ varies by six orders of
magnitude across the network.  Normalizing each bus's residual by its
local admittance or nominal power removes the $10^6$ conditioning that
caused the flat-PINN divergence in the LVN experiments:

\begin{equation}
  \widetilde{\Delta P}_i = \frac{\Delta P_i}{|S_i^{\rm set}| + \epsilon},
  \qquad
  \widetilde{\Delta Q}_i = \frac{\Delta Q_i}{|S_i^{\rm set}| + \epsilon}
\end{equation}

This is the \texttt{setpoint} normalization mode described in the physics
loss reference document.

%==============================================================================
\section{Recommended Reporting Table}
%==============================================================================

\begin{table}[h]
\centering
\caption{Recommended metric set for a GNN AC-PF surrogate.
  All mismatch quantities use the masked convention
  (slack excluded from $\Delta P$; slack + PV excluded from $\Delta Q$).}
\label{tab:report}
\begin{tabular}{llll}
\toprule
Metric & Unit & Statistics & Axis \\
\midrule
$|V|$ error & p.u. & MAE, max & Accuracy \\
$\theta$ error & deg & MAE, max & Accuracy \\
$|\Delta P|$ (PV+PQ) & p.u. / MW & mean, p95, max & Feasibility \\
$|\Delta Q|$ (PQ)    & p.u. / MVAr & mean, p95, max & Feasibility \\
Convergence rate & \% buses $< 10^{-3}$\,p.u. & — & Feasibility \\
NR polish steps & count & mean over test set & Practical value \\
\bottomrule
\end{tabular}
\end{table}

The \textbf{NR polish step count} is a practical bridge metric: starting
from the surrogate's predicted voltage, run pandapower until convergence
and count how many NR iterations were required.  Compare against a flat
start ($|V|=1$\,p.u., $\theta=0$).  This quantifies the surrogate's value
as a warm-start initializer even when its raw residual is above solver
tolerance.

\paragraph{Why this table?}
\begin{itemize}[noitemsep]
  \item $|V|$ and $\theta$ errors confirm the surrogate learned the
    correct operating point.
  \item Mean/p95/max $|\Delta P|$, $|\Delta Q|$ describe the distribution
    of KCL violations across buses, not just the worst case.
  \item The convergence rate maps raw residuals onto a single binary
    judgement that mirrors industrial practice.
  \item The NR polish step count makes the accuracy--speed trade-off
    quantitative: a surrogate that halves NR iterations is concretely
    useful even if it is not itself converged.
\end{itemize}

\end{document}
